\documentclass[xetex,aspectratio=43]{beamer}

\usepackage{res/lections}

\preamble

\title{ЦАП и АЦП}

\begin{document}

    \titleslide

    \tocslide


%==============================================

\section{ЦАП}

\subsection{Электронные преобразователи}

\begin{frame}{ЦАП R-2R}

\href{http://en.wikipedia.org/wiki/Resistor_ladder}{R-2R}

Достоинства:
\begin{itemize}
\item Работает за константу по времени
\item Устроен просто
\item Для $N$ уровней сложность схемы пропорциональна $\log_2 N$
\end{itemize}
\pause

Недостатки:
\begin{itemize}
\item Из-за необходимости в точных резисторах стоил дорого
\item Неравномерная нагрузка на резисторы и входные каналы
\end{itemize}
\pause
\begin{itemize}
\item Мощные точные резисторы --- практически оксюморон (из-за этого дороговизна)
\end{itemize}

Использовался, как дешёвый синтезатор звука: \href{https://en.wikipedia.org/wiki/Covox_Speech_Thing}{Covox Speech Thing},
поддерживался играми Wolfenstein 3D и Space Quest. Платформы Specturm,  БК и PC (через LPT).
\end{frame}

\begin{frame}{Широтно-импульсный модулятор}
\begin{columns}
\begin{column}{0.4\textwidth}
\includegraphics[height=0.9\textheight]{img/\jobname/wim}
\end{column}
\begin{column}{0.6\textwidth}
  \begin{itemize}
  \item Работает небыстро (пропорционально количеству значений)
  \item Хорошо масштабируем (меняется только разрядность счётчика)
  \item Схема сглаживания может вызывать «нелинейность», но её легче всего компенсировать цифровой калибровкой.
  \end{itemize}
\end{column}
\end{columns}
\end{frame}

\subsection{Пространственные манипуляторы}
\begin{frame}
  \begin{itemize}
  \item Высокой точности
    \begin{itemize}
    \item \href{https://en.wikipedia.org/wiki/Stepper_motor}{Шаговый двигатель}
    \end{itemize}
  \item Больших скорости и диапазона
    \begin{itemize}
    \item Лифт
    \item Подъёмный кран и даже строительный вертолёт
    \end{itemize}
  \item Гибридные
    \begin{itemize}
    \item Приблизительное позиционирование низкой точности, фиксация, затем точная подстройка, требующая обратной связи при помощи \hyperlink{adc}{АЦП} и
          \hyperlink{cybcontrollers}{кибернетических регуляторов}.
    \end{itemize}
  \end{itemize}
\end{frame}

\hypertarget{adc}{\section{АЦП}}

\begin{frame}{Что делает АЦП?}

\begin{itemize}
\item \defn{Квантование}{разбиение диапазона значений на дискретные уровни}
\item \defn{Дискретизация}{преобразования непрерывного аналогового сигнала в последовательность отдельных числовых отсчетов}
\end{itemize}

\end{frame}


\subsection{Электронные преобразователи}
\begin{frame}{Компаратор}
\begin{center}
\includegraphics[width=0.7\textwidth]{img/\jobname/comparaor}
\end{center}
\end{frame}

\begin{frame}{Прямой АЦП}
\begin{columns}
\begin{column}{0.6\textwidth}
\includegraphics[width=\textwidth]{img/\jobname/directadc}
\end{column}
\begin{column}{0.4\textwidth}
  \begin{itemize}
  \item Работает за линейное время (но быстро, аналоговая схема)
  \item Сложность схемы пропорциональна количеству позиций
  \end{itemize}
\end{column}
\end{columns}
\end{frame}

\begin{frame}{Компенсационный}
\begin{center}
\includegraphics[height=0.9\textheight]{img/\jobname/compenseadc}
\end{center}
Обычно подбирает значение бинарным поиском. Сложность логарифмическая.
\end{frame}

\begin{frame}{Широтно-импульсный демодулятор}
Можно измерять время разрядки конденсатора через резистор. Когда конденсатор разрядится ниже определённой разности потенциалов, должен остановиться счётчик.
\end{frame}

\subsection{Пространственные датчики}

\begin{frame}{Коды Грея и Де-Брёйна в датчиках координат}

\begin{block}{\href{https://en.wikipedia.org/wiki/Gray_code}{Код Грея}}
\begin{itemize}
\tightlist
\item Используются в механических и оптико-механических датчиках
\item Представляет собой обычные двоичные числа упорядоченные таким образом, что при переходе к следующему значению меняется только 1 бит. Код Грея устраняет неоднозначность при дребезге АЦП
\item Используется при параллельной пространственной маркировке
\end{itemize}
\end{block}

\begin{block}{\href{https://en.wikipedia.org/wiki/De_Bruijn_sequence}{Последовательность де Брёйна}}
\begin{itemize}
\tightlist
\item Последовательность битов $N$, любой отрезок которой длиной не менее $M$ битов уникален
\item Используется при последовательной пространственной маркировке, например в \emph{линейных энкодерах}
\end{itemize}
\end{block}

\end{frame}

\begin{frame}{Датчик угловой и линейной скоростей}
Используется, например, в механических и старых (требующих специального коврика) оптических «мышах».

\href{https://en.wikipedia.org/wiki/Computer_mouse\#Mechanical_mice}{Мышь}
\end{frame}

\begin{frame}{Интерферометрические: нониус и «родственники»}
  \begin{itemize}
  \item \href{https://en.wikipedia.org/wiki/Vernier_scale\#Functioning}{Штангенциркуль}
  \item \href{https://skbis.ru/catalog}{Линейные энкодеры с аналоговой составляющей}
  \end{itemize}
\end{frame}

\hypertarget{cybcontrollers}{\section{Кибернетические регуляторы}}

\begin{frame}{Регулятор}

\defn{Регулятор}{математическая модель учёта данных обратной связи,
позволяющая оптимизировать параметры функции с целью минимизации нежелательных
эффектов}

А какие могут быть нежелательные эффекты?

\pause

\begin{itemize}
\tightlist
\item
  Сильное запаздывание
\item
  Слишком резкое изменение функции
\item
  Колебания
\end{itemize}

А почему они нежелательные?

\pause

Представьте себе, что они проявляются:

\begin{itemize}
\tightlist
\item
  при управлении лифтом
\item
  при управлении отоплением
\end{itemize}


\end{frame}

\begin{frame}{Пример: ПИД-регулятор}

Один из способов улучшить ситуацию --- применить подходящий регулятор, наример, \href{https://en.wikipedia.org/wiki/PID_controller}{ПИД}

\end{frame}

\postamble

\end{document}
